\chapter{Signal Processing and Hardware Oracles}

\section{The Cryptographic Nervous System}
The Sovereign Stack is a unified cybernetic organism. If Layer 1 (Physical Foundations) is the thermodynamic metabolism governed by dissipative structures and entropy, then Layer 2 is its cryptographic nervous system. The sole purpose of Layer 2 is to sample the continuous, chaotic physics of the real world and perfectly translate them into deterministic state changes for the Layer 3 Ledger. 

\section{The Oracle Problem and Edge Signatures}
Layer 2 is the most vulnerable boundary in the Sovereign Stack. While Layer 3 is protected by rigorous cryptography and Layer 1 physics are strictly immutable (as defined by Prigogine's formulations of open systems), Layer 2 bridges the two. If a malicious actor can spoof sensor data (e.g., claiming a node generated 100 kWh of solar energy when it did not), the ledger will mint fraudulent Value Tokens, causing systemic economic collapse. This is the Oracle Problem.

To prevent oracle attacks, the Oasis architecture mandates that raw analog signals be converted to digital formats and cryptographically signed at the absolute edge—within the silicon of the sensor itself—using Hardware Security Modules (HSMs) or Secure Enclaves (e.g., ATECC608A).

\section{Scenario 0: Exergy Tracking and Server Telemetry}
Before simulating ecology or society, the node must monitor its own thermodynamic existence. Scenario 0 dictates that the computational cost of the software development and ledger consensus must be physically measured and priced in exact exergy.

\subsection{Current Shunt Monitoring and ADC Quantization}
Server power draw is continuously monitored using high-side current sense amplifiers (e.g., INA226) measuring the voltage drop across a precision shunt resistor ($R_{shunt}$). The true power $P$ is calculated as:
\begin{equation}
P = V_{bus} \times \left( \frac{V_{shunt}}{R_{shunt}} \right)
\end{equation}

Because the physical environment is noisy, the analog voltage must be digitized by an Analog-to-Digital Converter (ADC). The resolution of this conversion dictates the financial accuracy of the network. The theoretical Signal-to-Noise Ratio (SNR) due to quantization noise for an $N$-bit ADC is given by:
\begin{equation}
SNR = 6.02N + 1.76 \text{ [dB]}
\end{equation}
For precise exergy billing, a 16-bit ADC is utilized, providing $\approx 98$ dB of dynamic range \cite{kester2005data}. If a node operator attempts to spoof the voltage reading by injecting external analog noise, the digital filters (e.g., moving average FIR filters) running inside the secure enclave will flag the high-variance signal and suspend the node's minting privileges.

\section{Sensor Fusion and Dead Reckoning (Scenario Beta: Transit)}
In logistical scenarios (e.g., decentralized cargo bike networks), the node must track physical assets. However, relying purely on external GPS is a major vulnerability (susceptible to spoofing or urban canyon signal loss). 

To secure transit telemetry, Layer 2 employs Sensor Fusion, combining low-frequency GNSS data with high-frequency Inertial Measurement Unit (IMU) data. The IMU provides continuous acceleration ($a$) and angular velocity ($\omega$). A Kalman Filter is utilized to mathematically estimate the true physical state of the asset by minimizing the mean square error of the noisy sensor inputs:
\begin{equation}
\hat{x}_k = \hat{x}_{k}^{-} + K_k (z_k - H \hat{x}_{k}^{-})
\end{equation}
where $\hat{x}_k$ is the updated state estimate, $\hat{x}_{k}^{-}$ is the a priori state estimate, $K_k$ is the Kalman gain, $z_k$ is the actual measurement, and $H$ is the observation model mapping the true state space into the observed space. This recursive algorithm ensures that even if GPS is jammed, the asset's trajectory is cryptographically maintained via dead reckoning, preventing "ghost routing" attacks against the logistics ledger \cite{grewal2014kalman}.

\section{Structural Health and Vibration Analysis (Scenario Sigma)}
As physical infrastructure endures the friction of the legacy world, it mathematically accelerates toward entropy (Scenario Sigma: Burnout). Rather than relying on subjective human visual inspections, Layer 2 deploys piezoelectric accelerometers directly bonded to critical structural members (e.g., load-bearing beams, microgrid turbines).

The raw acoustic emission data is processed at the edge using the discrete Fast Fourier Transform (FFT) to convert the time-domain vibration signal $x(n)$ into the frequency domain $X(k)$:
\begin{equation}
X(k) = \sum_{n=0}^{N-1} x(n) e^{-j \frac{2\pi}{N} k n}
\end{equation}
As a structural material develops micro-fractures, its stiffness $k$ decreases, which physically shifts its resonant frequency downward ($\omega = \sqrt{k/m}$). By continuously running FFTs on the edge microcontrollers, the system can autonomously detect sub-millimeter structural degradation and proactively mint a \texttt{MaintenanceBounty} before a catastrophic kinetic failure occurs \cite{brincker2015introduction}.

\section{Cryptographic Proof of Labor (Scenarios Pi \& Kappa)}
To verify physical labor for guilds (Scenario Pi) and validate apprentice hours (Scenario Kappa) without requiring centralized human managers, the node tracks the telemetry of the physical tools themselves. High-current tools (welders, CNC spindles) are fitted with Hall-effect current sensors (e.g., ACS712) and 6-DoF IMUs.

By correlating the spatial orientation of the tool with the precise electrical current draw waveform, the edge node's embedded DSP generates a unique "labor signature." Attempting to spoof work by simply leaving a tool running idle will fail, as the lack of mechanical resistance drastically alters the phase angle between the current and voltage. This allows Layer 2 to automatically sign and emit a \texttt{ProofOfLabor} payload to the Layer 3 ledger, converting physical exertion directly into economic Value Tokens.

\vspace{1cm}
\begin{thebibliography}{99}
\bibitem{kester2005data}
Kester, W. (2005). \textit{Data conversion handbook}. Newnes.

\bibitem{grewal2014kalman}
Grewal, M. S., \& Andrews, A. P. (2014). \textit{Kalman filtering: Theory and Practice with MATLAB}. John Wiley \& Sons.

\bibitem{brincker2015introduction}
Brincker, R., \& Ventura, C. (2015). \textit{Introduction to operational modal analysis}. John Wiley \& Sons.
\end{thebibliography}
